% ============================================================
% Chapter 11: Non-Parametric Methods
% ============================================================
\chapter{Non-Parametric Methods}
\label{ch:nonparametric}

\begin{center}
\textit{``Do not make unnecessary assumptions.''} --- William of Ockham
\end{center}

\section{Motivating Example}
A psychologist compares well-being scores (ordinal, 1--10 scale) between meditation and control groups. Data are clearly non-normal. Non-parametric methods provide valid alternatives.

\section{Why Non-Parametric?}

\begin{itemize}
\item Ordinal or rank data.
\item Small samples where normality cannot be verified.
\item Heavily skewed distributions or outliers.
\item Unknown distribution: no parametric model desired.
\end{itemize}

\begin{warning}
Non-parametric tests are generally \textbf{less powerful} than parametric counterparts when parametric assumptions hold, but they are more \textbf{robust}.
\end{warning}

\section{Sign Test}
\index{sign test}

For $H_0:\text{median}=m_0$: count $S^+=\#\{i:X_i>m_0\}$. Under $H_0$: $S^+\sim\mathrm{Bin}(n,1/2)$.

\section{Wilcoxon Signed-Rank Test}
\index{Wilcoxon!signed-rank}

Compute $D_i=X_i-m_0$, rank $|D_i|$, and $W^+=\sum_{\{D_i>0\}} R_i$.

\section{Mann--Whitney / Wilcoxon Rank-Sum Test}
\index{Mann--Whitney}

Combine and rank both samples. $U_1=\sum R_i - n_1(n_1+1)/2$. For large samples:
\[
Z = \frac{U-n_1n_2/2}{\sqrt{n_1n_2(n_1+n_2+1)/12}} \xrightarrow{\mathcal{L}} \mathcal{N}(0,1).
\]

\section{Kruskal--Wallis Test}
\index{Kruskal--Wallis}

Extension to $k$ groups: $H = \frac{12}{n(n+1)}\sum R_j^2/n_j - 3(n+1) \sim \chi^2(k-1)$.

\section{Kolmogorov--Smirnov Test}
\index{Kolmogorov--Smirnov}

One-sample: $D_n=\sup_x|\hat{F}_n(x)-F_0(x)|$. Two-sample: $D_{n_1,n_2}=\sup_x|\hat{F}_{n_1}(x)-\hat{G}_{n_2}(x)|$.

\section{Kernel Density Estimation}
\index{KDE}\index{kernel density estimation}

\begin{definition}
$\hat{f}_h(x) = \frac{1}{nh}\sum_{i=1}^n K\!\left(\frac{x-X_i}{h}\right)$, where $K$ is a kernel and $h$ the bandwidth.
\end{definition}

\begin{intuition}
Small $h$: noisy (low bias, high variance). Large $h$: smooth (high bias, low variance). The bias-variance tradeoff in density estimation.
\end{intuition}

\section{Spearman Correlation}
\index{Spearman correlation}

$r_s = 1-\frac{6\sum d_i^2}{n(n^2-1)}$ where $d_i=\text{rank}(x_i)-\text{rank}(y_i)$. Measures monotonic (not necessarily linear) association.

\section{Python Implementation}

\begin{python}
\begin{minted}{python}
import numpy as np
from scipy import stats
import matplotlib.pyplot as plt

# Wilcoxon signed-rank
before = np.array([5.2, 4.8, 6.1, 5.5, 4.9, 5.8, 6.3, 5.1, 4.7, 5.4])
after = np.array([4.8, 4.2, 5.5, 5.0, 4.3, 5.1, 5.7, 4.5, 4.1, 4.9])
stat, p = stats.wilcoxon(before, after, alternative='greater')
print(f"Wilcoxon signed-rank: W={stat:.1f}, p={p:.4f}")

# Mann-Whitney
A = np.array([23, 25, 28, 30, 27, 24, 26, 29, 31, 22])
B = np.array([18, 20, 22, 19, 21, 17, 23, 20, 19, 18])
u, p_mw = stats.mannwhitneyu(A, B, alternative='two-sided')
print(f"Mann-Whitney: U={u:.1f}, p={p_mw:.4f}")

# Kruskal-Wallis
h, p_kw = stats.kruskal([6.4,7.1,6.8,7.3,6.9], [8.2,7.9,8.5,8.1,8.4], [5.5,6.2,5.8,6.0,5.7])
print(f"Kruskal-Wallis: H={h:.4f}, p={p_kw:.4f}")

# KS test
data = np.random.exponential(2, 50)
d, p_ks = stats.kstest(data, 'norm', args=(np.mean(data), np.std(data)))
print(f"KS (normality): D={d:.4f}, p={p_ks:.4f}")

# KDE
data_kde = np.concatenate([np.random.normal(-2, 0.8, 100), np.random.normal(2, 1.2, 150)])
x_grid = np.linspace(-6, 7, 300)
fig, ax = plt.subplots(figsize=(10, 5))
ax.hist(data_kde, bins=30, density=True, alpha=0.3)
for bw in [0.2, 0.5, 1.0]:
    kde = stats.gaussian_kde(data_kde, bw_method=bw)
    ax.plot(x_grid, kde(x_grid), lw=2, label=f'h={bw}')
ax.legend(); ax.set_title("Kernel Density Estimation")
plt.savefig("ch11_kde.pdf"); plt.show()

# Spearman
x = np.arange(1, 11); y = x**2
r_s, p_s = stats.spearmanr(x, y)
print(f"Spearman: r_s={r_s:.4f} (Pearson: {stats.pearsonr(x,y)[0]:.4f})")
\end{minted}
\end{python}

\section{Exercises}

\begin{exercise}[$\star$]
Two algorithms tested on 8 datasets. Is algorithm A significantly faster? (Mann--Whitney)
\end{exercise}

\begin{exercise}[$\star\star$]
Compare Type I error rates of $t$-test and Wilcoxon for log-normal data via simulation.
\end{exercise}

\begin{exercise}[$\star\star$]
Implement a Gaussian KDE. Study bandwidth effect on MISE for a Gaussian mixture.
\end{exercise}

\begin{exercise}[$\star\star\star$ -- Project]
Collect non-normal real data. Apply all non-parametric methods from this chapter. Compare with parametric results.
\end{exercise}

\begin{keyformulas}
\begin{align*}
U &= \sum R_i - n_1(n_1+1)/2 & H &= \frac{12}{n(n+1)}\sum R_j^2/n_j-3(n+1) \\
D_n &= \sup_x|\hat{F}_n(x)-F_0(x)| & r_s &= 1-6\sum d_i^2/[n(n^2-1)] \\
\hat{f}_h(x) &= \frac{1}{nh}\sum K\!\left(\frac{x-X_i}{h}\right)
\end{align*}
\end{keyformulas}
